\section*{ОБОЗНАЧЕНИЯ И СОКРАЩЕНИЯ}
\addcontentsline{toc}{section}{ОБОЗНАЧЕНИЯ И СОКРАЩЕНИЯ}

БД -- база данных.

ИС -- информационная система.

ИС~УП -- информационная система управления проектами.

ИТ -- информационные технологии.

ПО -- программное обеспечение.

РФ -- Российская Федерация.

СУБД -- система управления базами данных.

ТЗ -- техническое задание.

ТП -- технический проект.

API (Application Programming Interface) -- программный интерфейс приложения, набор соглашений и средств взаимодействия одной программной системы с другой.

CPM (Critical Path Method) -- метод критического пути, метод сетевого планирования работ проекта.

CRM (Customer Relationship Management) -- система управления взаимоотношениями с клиентами.

CSRF (Cross-Site Request Forgery) -- межсайтовая подделка запроса, разновидность атаки на веб-приложения, при которой запрос от имени пользователя выполняется со стороннего сайта.

CSS (Cascading Style Sheets) -- каскадные таблицы стилей, язык описания внешнего вида веб-страниц.

DRF (Django REST Framework) -- библиотека для построения программных интерфейсов в формате REST поверх фреймворка Django.

ER (Entity-Relationship) -- модель <<сущность -- связь>>, графический способ описания структуры данных.

ERP (Enterprise Resource Planning) -- система планирования ресурсов предприятия.

HTML (HyperText Markup Language) -- язык гипертекстовой разметки веб-страниц.

HTTP (HyperText Transfer Protocol) -- протокол прикладного уровня для передачи гипертекстовых документов.

HTTPS (HyperText Transfer Protocol Secure) -- расширение протокола HTTP с шифрованием передаваемых данных средствами TLS.

JSON (JavaScript Object Notation) -- текстовый формат обмена структурированными данными.

ORM (Object-Relational Mapping) -- объектно-реляционное отображение, технология связывания объектов программы со строками реляционной базы данных.

PBKDF2 (Password-Based Key Derivation Function 2) -- функция формирования криптографического ключа из пароля с применением соли и многократного хеширования.

PERT (Program Evaluation and Review Technique) -- техника анализа и оценки программ, метод сетевого планирования.

PMBOK (Project Management Body of Knowledge) -- свод знаний по управлению проектами, выпускаемый Project Management Institute.

PMI (Project Management Institute) -- Институт управления проектами, профессиональная организация в области проектного менеджмента.

REST (Representational State Transfer) -- архитектурный стиль построения распределенных систем, основанный на стандартных операциях протокола HTTP над ресурсами.

SaaS (Software as a Service) -- модель распространения программного обеспечения, при которой приложение размещается на инфраструктуре поставщика и предоставляется пользователям через сеть Интернет.

SAFe (Scaled Agile Framework) -- масштабированный гибкий фреймворк, набор принципов и практик для применения гибких методологий в крупных организациях с несколькими командами разработки.

SPA (Single Page Application) -- одностраничное веб-приложение.

SQL (Structured Query Language) -- язык структурированных запросов к реляционным базам данных.

UML (Unified Modeling Language) -- унифицированный язык моделирования, применяемый при разработке программного обеспечения.

URL (Uniform Resource Locator) -- единый указатель ресурса в сети Интернет.

UTF-8 (Unicode Transformation Format, 8-bit) -- кодировка символов Юникода переменной длины, обратно совместимая с ASCII.

XP (eXtreme Programming) -- экстремальное программирование, гибкая методология разработки программного обеспечения.

XSS (Cross-Site Scripting) -- межсайтовый скриптинг, разновидность атаки на веб-приложения через внедрение исполняемого кода в страницы, отображаемые другим пользователям.
